% =====================================================================
% Mirante dos Dados — Working Paper #8
% Custo real de operar uma plataforma de dados públicos no Databricks
% Free Edition: visibilidade fim-a-fim a partir de system tables
% (junho/2025 a abril/2026)
%
% Autor: Leonardo Chalhoub
% Versão: 1.0 — Abril de 2026
% Compilação: pdflatex articles/finops-plataforma-dados-publicos.tex (2x)
% =====================================================================

\documentclass[12pt,a4paper,oneside]{article}

\usepackage[T1]{fontenc}
\usepackage[utf8]{inputenc}
\usepackage[brazilian]{babel}
\usepackage{newtxtext}
\usepackage{newtxmath}
\usepackage[section]{placeins}
\usepackage[a4paper,top=3cm,bottom=2cm,left=3cm,right=2cm]{geometry}
\usepackage{setspace}
\usepackage{indentfirst}
\usepackage{titlesec}
\usepackage{caption}
\usepackage{booktabs}
\usepackage{array}
\usepackage{multirow}
\usepackage{graphicx}
\usepackage[colorlinks=true,
            urlcolor=blue!60!black,
            linkcolor=black,
            citecolor=black,
            breaklinks=true]{hyperref}
\usepackage{url}
\usepackage{xurl}
\Urlmuskip=0mu plus 1mu\relax
\usepackage{float}
\usepackage{tcolorbox}

\graphicspath{{figures-finops/}}

\onehalfspacing
\setlength{\parindent}{1.25cm}
\setlength{\parskip}{0pt}

\titleformat{\section}{\normalfont\bfseries\large\uppercase}{\thesection}{0.6em}{}
\titleformat{\subsection}{\normalfont\bfseries\normalsize}{\thesubsection}{0.6em}{}
\titlespacing*{\section}{0pt}{18pt}{8pt}
\titlespacing*{\subsection}{0pt}{12pt}{6pt}

\captionsetup[table]{labelfont=bf,name=Tabela,labelsep=endash,
                     position=above,justification=centering,font=small}
\captionsetup[figure]{labelfont=bf,name=Figura,labelsep=endash,
                      position=below,justification=centering,font=small}
\captionsetup{skip=6pt}

\newcommand{\us}{US\$\,}

% Stamp de compilação injetado pelo CI (deploy-pages.yml).
% Para builds locais, este arquivo default é usado — basta um pdflatex direto
% e a capa exibe "(dev local)". Em CI, este arquivo é sobrescrito antes da
% compilação com timestamp BRT + SHA do commit que disparou o deploy.
\newcommand{\COMPILESTAMP}{compila\c{c}\~ao local}
\newcommand{\COMPILESHA}{dev}


\begin{document}

% ---------- CAPA ---------------------------------------------------------
\begin{titlepage}
  \begin{center}
    \vspace*{1.5cm}
    {\small\textsc{Working Paper n. 8}}\\
    {\small Versão 1.0 \textemdash{} Abril de 2026}\\[2pt]
    {\footnotesize Compilada em \COMPILESTAMP\ \textbullet\ \texttt{\COMPILESHA}}

    \vfill

    {\LARGE\bfseries
     CUSTO REAL DE OPERAR UMA PLATAFORMA

     DE DADOS PÚBLICOS NO

     DATABRICKS FREE EDITION:

     VISIBILIDADE FIM-A-FIM A PARTIR DE

     \textit{SYSTEM TABLES} (JUNHO/2025\textendash ABRIL/2026)}

    \vspace{1cm}
    {\large\itshape
     Real cost of operating a public-data platform on Databricks
     Free Edition: end-to-end visibility from system tables
     (June 2025\textendash April 2026)}

    \vfill

    {\large\textbf{Leonardo Chalhoub}}\\[4pt]
    {\small Pesquisador independente \textemdash{} engenharia de dados,
     governança de custo (\textit{FinOps}), arquitetura
     \textit{lakehouse} e instrumentação de \textit{pipelines}.}\\[2pt]
    {\small\href{mailto:leonardochalhoub@gmail.com}{leonardochalhoub@gmail.com}}\\
    {\small\href{https://github.com/leonardochalhoub/mirante-dos-dados-br}{github.com/leonardochalhoub/mirante-dos-dados-br}}

    \vfill

    {\small Brasil \\ Abril de 2026}
  \end{center}
\end{titlepage}

% ---------- RESUMO -------------------------------------------------------
\section*{Resumo}
\addcontentsline{toc}{section}{Resumo}
\begin{singlespace}\small\noindent
A literatura sobre arquiteturas \textit{lakehouse} e plataformas
\textit{cloud-native} concentra-se em desempenho e governança;
estudos do \textit{custo real} de operar uma plataforma pequena, em
nível indivíduo \textit{single-developer}, são escassos. Este
\textit{working paper} apresenta auditoria fim-a-fim do gasto de
uma plataforma de dados públicos brasileiros, autônoma e
\textit{open-source}, hospedada no Databricks Free Edition,
abrangendo a janela completa entre 11 de junho de 2025 e 28 de
abril de 2026 (\textbf{322 dias contínuos, sem lacunas}).
\textbf{Materiais e métodos.} A fonte primária é constituída pelas
\textit{system tables} gerenciadas pelo Databricks
(\texttt{system.billing.usage}, \texttt{system.billing.list\_prices},
\texttt{system.lakeflow.job\_run\_timeline}). O custo de cada
registro de \textit{billing} é calculado como
\(\text{cost} = \text{usage\_quantity}_{\mathrm{DBU}} \times
\text{list\_price}_{\mathrm{USD/DBU}}\), onde \texttt{list\_price} é
versionado no tempo: para cada \textit{record} de uso, recupera-se a
linha de preço cuja janela
\([\text{price\_start\_time}, \text{price\_end\_time})\) contém
\texttt{usage\_start\_time}, casando \((\text{sku\_name},
\text{cloud}, \text{usage\_unit})\). O \textit{ledger} resultante
é materializado como duas tabelas Delta em camada \textit{silver}
(\texttt{finops\_run\_costs}, granularidade
\(\langle\text{job\_id, run\_id}\rangle\); \texttt{finops\_daily\_spend},
granularidade \(\langle\text{usage\_date, product, workload\_class}\rangle\))
e um par \textit{gold} agregado para consumo por \textit{front-end}.
\textbf{Achados.} O custo total da janela é \us 69{,}99,
distribuídos em 504,00 DBUs e 91 \textit{job runs} registradas. A
divisão por produto da Databricks revela \us 36{,}03 (51,5\,\%) em
\texttt{JOBS}, \us 16{,}33 (23,3\,\%) em \texttt{SQL warehouses},
\us 6{,}66 (9,5\,\%) em \texttt{NETWORKING}, \us 5{,}09 (7,3\,\%) em
\texttt{DEFAULT\_STORAGE}, \us 4{,}92 (7,0\,\%) em
clusters interativos, \us 0{,}82 (1,2\,\%) em \textit{predictive
optimization} e \us 0{,}14 (0,2\,\%) em DLT.
\textbf{O achado-chave é a fração de \textit{wasted spend}}: do
gasto atribuível a \textit{job runs} (\us 36{,}03), apenas
\us 16{,}62 (46,1\,\%) corresponde a execuções bem-sucedidas; os
demais \textbf{\us 19{,}41 (53,9\,\%)} foram consumidos em
\textit{runs} que falharam (30 execuções, \us 11{,}72) ou foram
canceladas (30 execuções, \us 7{,}69). DBUs continuam sendo
faturadas até o crash ou o cancelamento manual. O \textit{run rate}
mensal de \textit{storage} (\texttt{DEFAULT\_STORAGE}), pro-rateado
pelos 19 dias com cobrança ativa, é \us 8{,}01/mês
(\us 0{,}267/dia), com projeção anual de \us 97{,}40. Pro-rateando
esse \textit{storage} pela participação em \textit{bytes} de cada
\textit{vertical} em \textit{bronze}, a maior fatia recai sobre o
\textit{vertical} \textit{Bolsa Família} (~99\,\% dos \textit{bytes}
de \textit{bronze}, \(\sim\)\us 7{,}90/mês de \textit{storage}
imputado). \textbf{Discussão e contribuição.} O \textit{paper}
documenta um padrão arquitetural simples (\textit{system tables}
como \textit{bronze} para a vertical FinOps, derivada por
\textit{batch overwrite} \textit{silver}/\textit{gold}, exposta via
JSON estático ao \textit{front-end}) que torna a auditoria de custo
ao nível indivíduo um subproduto do \textit{lakehouse}, não uma
ferramenta separada. A magnitude do \textit{wasted spend}
(\(>\!50\,\%\) do gasto em \textit{jobs}) reproduz, em escala
individual, padrões reportados pela literatura cinza FinOps em
ambientes empresariais maiores, sugerindo que o ciclo
\textit{visibility \textrightarrow{} allocation \textrightarrow{}
optimization} é estruturalmente independente de escala. O
\textit{pipeline} é \textit{open-source} e o \textit{ledger}
versionado é publicado junto com este \textit{working paper}.

\vspace{0.4cm}\noindent\textbf{Palavras-chave:}
governança de custo em nuvem; \textit{FinOps}; \textit{lakehouse};
Databricks; \textit{system tables}; \textit{wasted spend};
arquitetura \textit{medallion}; observabilidade de \textit{pipelines};
Delta Lake; \textit{billing-driven engineering}.
\end{singlespace}

% ---------- ABSTRACT -----------------------------------------------------
\section*{Abstract}
\addcontentsline{toc}{section}{Abstract}
\begin{singlespace}\small\noindent
The literature on \textit{lakehouse} architectures and cloud-native
data platforms focuses on performance and governance; studies of the
\textit{actual cost} of operating a small, single-developer platform
are scarce. This \textit{working paper} presents an end-to-end audit
of the spend of an open-source Brazilian public-data platform
hosted on Databricks Free Edition, covering the full window between
11 June 2025 and 28 April 2026 (\textbf{322 contiguous days,
no gaps}). \textbf{Methods.} The primary data source is the
Databricks-managed \textit{system tables}
(\texttt{system.billing.usage}, \texttt{system.billing.list\_prices},
\texttt{system.lakeflow.job\_run\_timeline}). Per-record cost is
\(\text{cost} = \text{usage\_quantity}_{\mathrm{DBU}} \times
\text{list\_price}_{\mathrm{USD/DBU}}\) with time-versioned price
joined on the interval that contains \texttt{usage\_start\_time}
(matching SKU, cloud and unit). The resulting ledger is materialized
in two Delta \textit{silver} tables and two \textit{gold}
aggregates feeding a static JSON consumed by the front-end.
\textbf{Findings.} Total window spend is \us 69.99 over 504.00 DBUs
and 91 job runs. Product split: \us 36.03 \texttt{JOBS} (51.5\,\%),
\us 16.33 \texttt{SQL} warehouses, \us 6.66 \texttt{NETWORKING},
\us 5.09 \texttt{DEFAULT\_STORAGE}, \us 4.92 interactive clusters,
\us 0.82 \textit{predictive optimization}, \us 0.14 DLT. Headline
finding: of the \textit{job} spend (\us 36.03), only \us 16.62
(46.1\,\%) corresponds to successful runs; \textbf{\us 19.41
(53.9\,\%)} was consumed by failed (30 runs, \us 11.72) or cancelled
(30 runs, \us 7.69) executions, since DBUs are billed until the
crash or manual cancellation. Storage run rate is \us 8.01/month
(\us 0.267/day; \us 97.40/year projection). Pro-rating storage by
each vertical's share of \textit{bronze bytes}, the
\textit{Bolsa Família} vertical absorbs almost the entire share
($\sim$99\,\%, $\sim$\us 7.90/month). \textbf{Discussion.} The
paper documents a simple architectural pattern (system tables as
bronze for a FinOps vertical, batch-overwritten silver/gold, served
via static JSON to the front) that makes individual-level cost
auditing a byproduct of the lakehouse rather than a separate tool.
The magnitude of wasted spend ($>$50\,\% of job spend) mirrors,
at individual scale, FinOps grey-literature findings in larger
enterprise contexts, suggesting that the
\textit{visibility \textrightarrow{} allocation \textrightarrow{}
optimization} cycle is structurally scale-invariant. The pipeline
is open-source and the versioned ledger is published with this
paper.

\vspace{0.4cm}\noindent\textbf{Keywords:}
cloud cost governance; FinOps; lakehouse; Databricks; system tables;
wasted spend; medallion architecture; pipeline observability;
Delta Lake; billing-driven engineering.
\end{singlespace}

\newpage

% ---------- 1. INTRODUÇÃO ------------------------------------------------
\section{Introdução}

A literatura técnica sobre arquiteturas \textit{lakehouse} é
abundante em três frentes: \textit{benchmarks} de desempenho entre
formatos abertos como Delta Lake, Apache Iceberg e Apache Hudi
(ARMBRUST et al., 2021); padrões de governança via Unity Catalog,
\textit{table-level RBAC} e \textit{data contracts}
(DATABRICKS, 2024); e padrões arquiteturais como o \textit{medallion}
(\textit{bronze} \textrightarrow{} \textit{silver}
\textrightarrow{} \textit{gold}) e \textit{data mesh}. Entretanto,
o estudo do \textit{custo monetário real} de operar uma plataforma
de pequeno porte permanece sub-documentado, especialmente em
ambientes de desenvolvimento individual ou de pesquisa
\textit{single-developer}.

A literatura cinza de \textit{FinOps} \textemdash{} originada na
prática profissional do FinOps Foundation (FINOPS FOUNDATION, 2023)
\textemdash{} reporta sistematicamente que organizações que adotam
plataformas \textit{cloud-native} sem governança específica de
custo desperdiçam entre 30\,\% e 35\,\% do gasto total
(FLEXERA, 2024). Esses números são produzidos em ambientes
empresariais com dezenas a milhares de \textit{users}, e geralmente
são derivados de pesquisas auto-reportadas ou de auditorias internas
não-públicas. Resta em aberto: o mesmo padrão emerge em escala
individual? Um \textit{single developer} responsável pela ingestão,
transformação e \textit{export} de microdados públicos brasileiros
gera \textit{wasted spend} similar?

Este \textit{working paper} responde com auditoria fim-a-fim de
uma plataforma específica: o conjunto de \textit{pipelines}
\textit{open-source} responsável pela transformação reproduzível
de microdados públicos brasileiros (Bolsa Família via Portal da
Transparência da CGU; equipamentos médicos via DATASUS/CNES;
emendas parlamentares via CGU; vínculos formais via PDET/MTE; AIH
hospitalares via DATASUS/SIH-RD) hospedados no plano gratuito do
Databricks (\textit{Free Edition}). A janela de observação cobre
\textbf{322 dias contínuos, sem lacunas} \textemdash{} 100\,\% do
histórico operacional desde a primeira execução de \textit{job}
(11 de junho de 2025) até a data de compilação deste artigo
(28 de abril de 2026).

A contribuição metodológica é a demonstração de que o
\textit{Single Source of Truth} de custo \textemdash{} as
\textit{system tables} mantidas pelo próprio provedor cloud
\textemdash{} pode ser materializado como \textit{vertical}
adicional dentro da mesma arquitetura \textit{medallion} já
operacional, com \textit{bronze} ``virtual'' (referência direta à
\texttt{system.billing.usage}), \textit{silver} de granularidade
fina (uma linha por \textit{job run}) e \textit{gold} pivotada para
consumo via \textit{front-end} estático. A auditoria de custo
torna-se subproduto natural da \textit{platform engineering},
não ferramenta externa requerendo \textit{agents},
\textit{exporters} ou \textit{scrape-and-store} de APIs proprietárias.

A contribuição empírica é o relato detalhado, em granularidade de
\textit{job-run} individual, do gasto desta plataforma específica.
O resultado-chave \textemdash{} \textbf{53,9\,\% do gasto em
\textit{jobs} foi consumido por \textit{runs} que falharam ou foram
canceladas} \textemdash{} é a primeira contribuição quantitativa
deste \textit{paper}.

\subsection*{Estrutura do artigo}

A Seção~\ref{sec:methods} descreve a fonte de dados (\textit{system
tables}), o algoritmo de \textit{pricing} versionado no tempo, a
arquitetura \textit{silver}/\textit{gold} e o processo de
\textit{export} JSON. A Seção~\ref{sec:results} apresenta os
achados em sete subseções: panorama agregado, série temporal
diária, \textit{breakdown} por produto, análise de
\textit{wasted spend}, custo por \textit{job}, perfil das
\textit{runs} mais caras, e \textit{run rate} de \textit{storage}
com pro-rate por \textit{vertical}. A Seção~\ref{sec:discussion}
discute o padrão arquitetural emergente
(\textit{billing-driven engineering}), as limitações da auditoria
e a implicação geral. A Seção~\ref{sec:limits} consolida limitações
metodológicas. A Seção~\ref{sec:reproducibility} fecha com a
\textit{checklist} de reprodutibilidade. A bibliografia adota o
padrão ABNT NBR~6023, estendido com URLs absolutas e \textit{access
timestamps} para todas as fontes \textit{web}.

% ---------- 2. MATERIAIS E MÉTODOS --------------------------------------
\section{Materiais e métodos}\label{sec:methods}

\subsection{Fonte de dados: \textit{system tables} do Databricks}

A Databricks expõe, no esquema \texttt{system}, três famílias de
\textit{system tables} relevantes a esta auditoria:
\begin{itemize}\setlength{\itemsep}{2pt}
  \item \texttt{system.billing.usage} \textemdash{} um \textit{record}
        por par \(\langle\text{recurso, hora}\rangle\), com colunas
        \texttt{usage\_quantity} (em DBUs), \texttt{usage\_unit},
        \texttt{usage\_start\_time}, \texttt{usage\_end\_time},
        \texttt{sku\_name}, \texttt{cloud},
        \texttt{usage\_metadata} (\textit{struct} com
        \texttt{job\_id}, \texttt{job\_run\_id}, \texttt{job\_name},
        \texttt{cluster\_id}, etc.) e
        \texttt{product\_features} (\textit{struct} com
        \texttt{is\_serverless}, \texttt{is\_photon}, etc.).
  \item \texttt{system.billing.list\_prices} \textemdash{} preços
        unitários (USD por DBU) versionados no tempo, com
        \texttt{price\_start\_time} e \texttt{price\_end\_time}
        (este último \texttt{NULL} para o preço corrente);
  \item \texttt{system.lakeflow.job\_run\_timeline} \textemdash{}
        uma linha por \textit{job run} (incluindo \textit{state
        transitions}), com \texttt{result\_state}
        (\texttt{SUCCEEDED}, \texttt{FAILED} \textemdash{} aliás
        \texttt{ERROR} no \textit{billing} \textemdash{},
        \texttt{CANCELLED}, etc.) e \texttt{period\_start\_time} /
        \texttt{period\_end\_time}.
\end{itemize}

Estas tabelas são gerenciadas pelo provedor: novas linhas são
inseridas continuamente com latência aproximada de uma hora para
\texttt{system.billing.usage}; \texttt{job\_run\_timeline} é
quase em tempo real após a \textit{run} terminar. O acesso é
\textit{read-only}.

\subsection{Algoritmo de \textit{pricing} versionado no tempo}

Para cada registro \(r\) de \texttt{system.billing.usage}, o custo
em USD é calculado como
\[
  \text{cost}(r)
  \;=\;
  \text{usage\_quantity}(r) \;\times\; \pi\!\bigl(\text{sku\_name}(r),
  \text{cloud}(r), \text{usage\_unit}(r),
  \text{usage\_start\_time}(r)\bigr),
\]
onde \(\pi(\cdot)\) é a função que recupera o preço unitário válido
no instante \texttt{usage\_start\_time}(r) a partir de
\texttt{system.billing.list\_prices}, condicionado a igualdade dos
três identificadores de SKU. Operacionalmente, isso corresponde ao
\textit{join} relacional
\begin{tcolorbox}[colback=gray!5,colframe=gray!40,boxrule=0.4pt,
                  arc=2pt,left=8pt,right=8pt,top=4pt,bottom=4pt]
\small\ttfamily
LEFT JOIN system.billing.list\_prices lp \\
\hspace*{1.5em}ON \ lp.sku\_name \ = u.sku\_name \\
\hspace*{2.5em}AND lp.cloud \ \ = u.cloud \\
\hspace*{2.5em}AND lp.usage\_unit = u.usage\_unit \\
\hspace*{2.5em}AND u.usage\_start\_time \(\geq\) lp.price\_start\_time \\
\hspace*{2.5em}AND (lp.price\_end\_time IS NULL \\
\hspace*{4em}OR \ u.usage\_start\_time \(<\) lp.price\_end\_time)
\end{tcolorbox}

\noindent A semântica de intervalo é \([\text{start}, \text{end})\):
o limite superior é exclusivo. O preço \texttt{NULL} no extremo
superior indica preço corrente sem revisão posterior; o predicado
\(\text{is null} \;\text{OR}\; r.\text{start} < \text{end}\) cobre
ambos os casos.

\subsection{Arquitetura \textit{silver}/\textit{gold}}

A vertical FinOps materializa quatro tabelas Delta no
\textit{catálogo} \texttt{mirante\_prd}:

\begin{itemize}\setlength{\itemsep}{4pt}
  \item \texttt{silver.finops\_run\_costs} \textemdash{} grão
        \(\langle\text{job\_id, run\_id}\rangle\). Agrega \textit{records}
        de \texttt{system.billing.usage} filtrados por
        \texttt{job\_run\_id} \texttt{IS NOT NULL}, junta com
        \texttt{system.lakeflow.job\_run\_timeline} para recuperar
        \texttt{result\_state}, \texttt{period\_start\_time} e
        \texttt{period\_end\_time}. Colunas: \texttt{job\_id},
        \texttt{run\_id}, \texttt{job\_name}, \texttt{run\_name},
        \texttt{result\_state}, \texttt{started\_at},
        \texttt{ended\_at}, \texttt{billed\_minutes},
        \texttt{dbus}, \texttt{cost\_usd}, \texttt{primary\_sku},
        \texttt{is\_serverless}, \texttt{is\_photon}, \texttt{day}.

  \item \texttt{silver.finops\_daily\_spend} \textemdash{} grão
        \(\langle\text{usage\_date, product, workload\_class,
        is\_serverless}\rangle\). Inclui \textbf{todos} os
        \textit{records} de \textit{billing} (não apenas os
        tagueados a \textit{job runs}), classificando cada produto
        em \textit{chargeable} (\texttt{JOBS}, \texttt{SQL},
        \texttt{INTERACTIVE}, \texttt{DLT}) ou \textit{overhead}
        (\texttt{NETWORKING}, \texttt{DEFAULT\_STORAGE},
        \texttt{PREDICTIVE\_OPTIMIZATION}). A distinção é
        operacional: \textit{chargeable} é diretamente atribuível
        a um \textit{job} ou \textit{warehouse}; \textit{overhead}
        é custo de plataforma sem \texttt{run\_id}, geralmente
        rateado proporcionalmente entre consumidores.

  \item \texttt{gold.finops\_run\_costs} \textemdash{} cópia
        enriquecida do \textit{silver} de mesmo nome. Adiciona
        \texttt{job\_name\_canonical} (com remoção do prefixo
        \texttt{[dev <user>]} para agrupar \textit{dev}/\textit{prod}
        \textit{runs} do mesmo \textit{job}), \texttt{is\_wasted}
        (\texttt{TRUE} se \texttt{result\_state} \(\in\)
        \(\{\texttt{ERROR}, \texttt{CANCELLED}\}\)) e
        \texttt{cost\_bucket} (categoriza em
        \texttt{micro} \(<\) \us 0{,}10, \texttt{small} \(<\)
        \us 0{,}50, \texttt{medium} \(<\) \us 2{,}00,
        \texttt{large} \(\geq\) \us 2{,}00).

  \item \texttt{gold.finops\_daily\_spend} \textemdash{} série
        diária pivoteada por produto, com colunas
        \texttt{cost\_jobs}, \texttt{cost\_sql},
        \texttt{cost\_interactive}, \texttt{cost\_dlt},
        \texttt{cost\_networking}, \texttt{cost\_storage},
        \texttt{cost\_pred\_opt}; agrega
        \texttt{cost\_chargeable\_total},
        \texttt{cost\_overhead\_total}, \texttt{cost\_total}; e
        deriva \texttt{cost\_total\_cumulative} (soma
        \textit{running} ao longo do histórico, via \textit{window
        function} \texttt{ROWS BETWEEN UNBOUNDED PRECEDING AND
        CURRENT ROW}).
\end{itemize}

A política de escrita é \texttt{overwrite} idempotente: cada
execução do \textit{job} \texttt{job\_finops\_refresh} reconstrói
as quatro tabelas a partir do estado corrente das \textit{system
tables}. Não há \textit{state} acumulado entre execuções; a única
informação histórica é o próprio \texttt{system.billing.usage},
mantido pelo provedor.

\subsection{\textit{Export} JSON e consumo \textit{front-end}}

Um \textit{notebook} de \textit{export}
(\texttt{export.finops\_summary\_json}) lê as duas tabelas
\textit{gold}, computa KPIs derivados em \textit{pandas}
(\textit{wasted spend} por desfecho, \textit{run rate} de
\textit{storage}, distribuição por produto, \textit{top-N} \textit{jobs}
e \textit{runs}, percentil 95 do custo por \textit{run}) e grava um
único objeto JSON em volume gerenciado, copiado para o repositório
\textit{open-source} via \textit{job} de sincronização. O
\textit{front-end} estático (React + Vite) renderiza a página
FinOps consumindo apenas esse JSON; nenhum acesso ao Databricks
ocorre em tempo de leitura pelo público.

\subsection{Janela de observação e completude}

A janela de observação cobre 11 de junho de 2025 a 28 de abril de
2026 \textemdash{} \(\mathbf{n = 322}\) dias contínuos. O dia
inicial (\(d_0 = 11/06/2025\)) corresponde ao primeiro registro
em \texttt{system.billing.usage} para o \textit{workspace} desta
plataforma, registrando custo em \texttt{DEFAULT\_STORAGE}. A
plataforma só passou a ter execuções de \textit{job} a partir de
abril de 2026, o que explica a concentração temporal extrema
discutida na Seção~\ref{sec:results}.

% ---------- 3. RESULTADOS -----------------------------------------------
\section{Resultados}\label{sec:results}

\subsection{Panorama agregado}

Na janela completa de 322 dias, foram registradas 91 execuções de
\textit{job} consumindo 503,99 DBUs, totalizando \us 69{,}99
em custo lifetime (Tabela~\ref{tab:overview}). A concentração
temporal é extrema: a quase totalidade do gasto ocorreu nos
\(\sim\)20 dias do mês de abril de 2026 (\us 69{,}48 nos últimos
30 dias; \us 64{,}15 nos últimos 7 dias \textemdash{} 91,7\,\% e
84,7\,\% do total \textit{lifetime}, respectivamente). A janela
completa antes desse pico observa apenas custos contínuos de
\textit{storage} e tráfego ocasional de \textit{warehouses}.

\begin{table}[H]
\centering
\caption{\textbf{Panorama agregado da plataforma}, 11/06/2025
        \textendash{} 28/04/2026 (\(n = 322\) dias).}
\label{tab:overview}
\begin{tabular}{lr}
\toprule
\textbf{Métrica} & \textbf{Valor} \\
\midrule
Custo total \textit{lifetime}                         & \us 69{,}99 \\
Custo últimos 30 dias                                 & \us 69{,}48 \\
Custo últimos 7 dias                                  & \us 64{,}15 \\
DBUs consumidas (\textit{lifetime})                   & 503,99 \\
Job \textit{runs} registradas                         & 91 \\
Custo médio por \textit{run}                          & \us 0{,}40 \\
Percentil 95 do custo por \textit{run}                & \us 1{,}49 \\
\textit{Run} mais cara registrada                     & \us 3{,}01 \\
Participação \textit{chargeable} no total             & 82,0\,\% \\
Participação \textit{overhead} no total               & 18,0\,\% \\
\textit{Wasted spend} (\textit{ERROR + CANCELLED})    & \us 19{,}41 \\
\textit{Wasted} \% do gasto em \textit{jobs}          & 53,9\,\% \\
\bottomrule
\end{tabular}
\end{table}

\subsection{\textit{Breakdown} por produto}

A Tabela~\ref{tab:byproduct} apresenta a decomposição por produto
da Databricks. Quatro padrões emergem.

\begin{table}[H]
\centering
\caption{\textbf{Custo \textit{lifetime} por produto Databricks}.
        \textit{Workload-class} divide custo em
        \textit{chargeable} (código do usuário) vs.\ \textit{overhead}
        (plataforma).}
\label{tab:byproduct}
\begin{tabular}{lcrr}
\toprule
\textbf{Produto} & \textbf{Classe} & \textbf{Custo (USD)} & \textbf{Share} \\
\midrule
\texttt{JOBS}                       & \textit{chargeable} & 36,03 & 51,5\,\% \\
\texttt{SQL}                        & \textit{chargeable} & 16,33 & 23,3\,\% \\
\texttt{NETWORKING}                 & \textit{overhead}   & \phantom{0}6,66 & \phantom{0}9,5\,\% \\
\texttt{DEFAULT\_STORAGE}           & \textit{overhead}   & \phantom{0}5,09 & \phantom{0}7,3\,\% \\
\texttt{INTERACTIVE}                & \textit{chargeable} & \phantom{0}4,92 & \phantom{0}7,0\,\% \\
\texttt{PREDICTIVE\_OPTIMIZATION}   & \textit{overhead}   & \phantom{0}0,82 & \phantom{0}1,2\,\% \\
\texttt{DLT}                        & \textit{chargeable} & \phantom{0}0,14 & \phantom{0}0,2\,\% \\
\midrule
\textbf{Total}                      & \textemdash{}       & \textbf{69,99} & \textbf{100,0\,\%} \\
\bottomrule
\end{tabular}
\end{table}

Primeiro, \texttt{JOBS} domina o gasto (51,5\,\%), refletindo a
arquitetura \textit{batch overwrite} adotada nos cinco
\textit{verticais} de microdados públicos. Segundo, \texttt{SQL
warehouses} contribui 23,3\,\% mesmo com uso ad-hoc, indicando que
operações de exploração (\texttt{DESCRIBE DETAIL},
\texttt{SHOW HISTORY}, \textit{queries} de validação) acumulam
custo não-trivial mesmo com \textit{auto-stop} agressivo (10
minutos). Terceiro, \texttt{INTERACTIVE} (\textit{notebooks}
ad-hoc fora de \textit{jobs}) responde por 7,0\,\%
\textemdash{} explorações manuais para investigar \textit{jobs}
quebrados. Quarto, a soma de \textit{overhead}
(\texttt{NETWORKING} \(+\) \texttt{DEFAULT\_STORAGE} \(+\)
\texttt{PREDICTIVE\_OPTIMIZATION}) atinge 18,0\,\% do total,
caracterizando custo estrutural da plataforma independente da
intensidade de uso.

\subsection{Análise do desfecho das \textit{runs}: \textit{wasted spend}}

A Tabela~\ref{tab:byoutcome} é a contribuição central deste
\textit{paper}. Das 91 \textit{job runs} registradas, 31 (34,1\,\%)
foram bem-sucedidas, 30 (33,0\,\%) falharam (\texttt{ERROR}) e 30
(33,0\,\%) foram canceladas (\texttt{CANCELLED}). A divisão de custo
não respeita essa proporção: \textit{runs} bem-sucedidas
consumiram apenas 46,1\,\% do gasto em \textit{jobs}, enquanto
\textit{runs} falhas e canceladas consumiram 32,5\,\% e 21,3\,\%,
respectivamente \textemdash{} \textbf{53,9\,\% de \textit{wasted
spend}}.

\begin{table}[H]
\centering
\caption{\textbf{Custo \textit{lifetime} por desfecho de \textit{run}}.
        Custo médio por \textit{run} e duração média também listados.
        Total de \textit{job} \textit{runs}: 91.}
\label{tab:byoutcome}
\begin{tabular}{lcrrrr}
\toprule
\textbf{Desfecho}        & \textbf{N} & \textbf{Custo (USD)} & \textbf{Share \%} & \textbf{Médio (USD)} & \textbf{Min/run} \\
\midrule
\texttt{SUCCEEDED}       & 31         & 16,62                & 46,1              & 0,54                 & 21,8 \\
\texttt{ERROR}           & 30         & 11,72                & 32,5              & 0,39                 & 22,9 \\
\texttt{CANCELLED}       & 30         & \phantom{0}7,69      & 21,3              & 0,26                 & 50,6 \\
\midrule
\textbf{Total}           & \textbf{91} & \textbf{36,03}      & \textbf{100,0}    & \textbf{0,40}        & \textbf{31,8} \\
\bottomrule
\end{tabular}
\end{table}

A duração média da \textit{run} cancelada (50,6 minutos) é
2,3$\times$ a da \textit{run} bem-sucedida (21,8 minutos), o que
sugere que cancelamentos ocorreram em \textit{jobs} que travaram
ou se tornaram intoleravelmente lentos, não em decisões precoces
de abortar. \textit{Runs} falhas têm duração média próxima da
mediana de execução com sucesso (22,9 vs.\ 21,8 minutos), o que é
consistente com falhas em \textit{tasks} próximos do fim do DAG
\textemdash{} um \textit{export} de JSON falhando após \textit{silver}
e \textit{gold} terem rodado, por exemplo. Cada hipótese é
verificável diretamente no \textit{ledger} \textit{silver}; este
\textit{paper} se limita a documentar a magnitude do desperdício.

A leitura econômica é direta: \textbf{quase 54\,\% de cada dólar
gasto em \textit{jobs} resultou em DBUs queimadas sem entregar valor
de \textit{output}}. Em escala individual, isso equivale a
\us 19{,}41; em escala empresarial onde gasto em \textit{jobs}
atinge dezenas de milhares de dólares ao mês, a magnitude do
\textit{wasted spend} sob a mesma proporção seria substancial.

\subsection{Custo por \textit{job}}

A Tabela~\ref{tab:byjob} apresenta os cinco \textit{jobs} mais
caros do \textit{lifetime}, agrupando \textit{dev}/\textit{prod}
\textit{runs} via \texttt{job\_name\_canonical}.

\begin{table}[H]
\centering
\caption{\textbf{Top-5 \textit{jobs} por custo \textit{lifetime}}.
        \textit{Wasted} \% é a fração do custo do próprio \textit{job}
        que recai em \textit{runs} \texttt{ERROR}/\texttt{CANCELLED}.}
\label{tab:byjob}
\small
\begin{tabular}{lrrrrr}
\toprule
\textbf{Job}                            & \textbf{Runs} & \textbf{Custo (USD)} & \textbf{Médio (USD)} & \textbf{Wasted (USD)} & \textbf{Wasted \%} \\
\midrule
refresh PBF                             & 8             & 9,68                 & 1,21                  & 4,75                  & 49,1 \\
\textit{manual silver}\(\to\)\textit{gold equipamentos} & 8 & 8,70    & 1,09                  & \textemdash{}          & \textemdash{} \\
refresh RAIS                            & 30            & 6,57                 & 0,22                  & \textemdash{}          & \textemdash{} \\
refresh Equipamentos                    & 13            & 5,01                 & 0,39                  & \textemdash{}          & \textemdash{} \\
refresh UroPro                          & \phantom{0}5   & 2,69                & 0,54                  & \textemdash{}          & \textemdash{} \\
\bottomrule
\end{tabular}
\end{table}

O \textit{job} \texttt{refresh PBF} concentra simultaneamente o
maior custo \textit{lifetime} (\us 9{,}68), o maior custo médio por
\textit{run} (\us 1{,}21) e o maior custo absoluto desperdiçado
(\us 4{,}75 em três \textit{runs} \texttt{ERROR}). Esse padrão
reflete a arquitetura específica do \textit{vertical} Bolsa Família
\textemdash{} o \textit{bronze} concentra 2,5 bilhões de linhas e
\(\sim\)42 GB de Delta storage; cada execução do DAG completo
re-processa essa massa.

\subsection{\textit{Top runs} mais caras}

A \textit{run} mais cara do \textit{lifetime} consumiu \us 3{,}01
em 65,9 minutos (8,59 DBUs), tendo desfecho \texttt{ERROR}, em
24/04/2026, no \textit{job} \texttt{refresh PBF}. Esse caso
isolado representa 4,3\,\% do gasto \textit{lifetime} da plataforma
inteira em uma única \textit{run} fracassada. As demais 24 \textit{runs}
do \textit{top-25} apresentam custos entre \us 0{,}63 e \us 1{,}77
\textemdash{} a distribuição é fortemente assimétrica à direita,
com cauda longa habitada por \textit{runs} grandes que falham.

\subsection{\textit{Storage}: custo contínuo, atribuição por \textit{vertical}}

A componente \texttt{DEFAULT\_STORAGE} é qualitativamente diferente
das demais: cobra continuamente, mesmo na ausência de execução de
\textit{jobs}. A Tabela~\ref{tab:storage} apresenta os indicadores
de \textit{run rate}.

\begin{table}[H]
\centering
\caption{\textbf{Indicadores de \textit{run rate} de \textit{storage}}.
        Pro-rate diário considera apenas dias com cobrança
        ($n = 19$).}
\label{tab:storage}
\begin{tabular}{lr}
\toprule
\textbf{Métrica}                                   & \textbf{Valor} \\
\midrule
Total \textit{lifetime}                            & \us 5{,}09 \\
Dias com cobrança                                  & 19 \\
Custo médio por dia (lifetime)                     & \us 0{,}268 \\
Custo médio por dia (últimos 30 dias)              & \us 0{,}267 \\
Run rate mensal (\(\times 30\))                    & \us 8{,}01 \\
Run rate anual (\(\times 365\))                    & \us 97{,}40 \\
Share de \textit{storage} no total \textit{lifetime} & 7,3\,\% \\
\bottomrule
\end{tabular}
\end{table}

Pro-rateando esse \textit{storage} pela participação em
\textit{bytes} de cada \textit{vertical} em \textit{bronze}
(Tabela~\ref{tab:storagepervert}), a maior parcela imputada recai
sobre Bolsa Família (98,9\,\% dos bytes; \(\sim\)\us 7{,}92/mês),
seguida por Incontinência Urinária (1,1\,\%, \(\sim\)\us 0{,}09/mês)
e os demais com fração inferior a 0,1\,\%. A pro-rate é
aproximação \textemdash{} a Databricks não imputa
\texttt{DEFAULT\_STORAGE} por \textit{schema} no \textit{billing};
trata-se de proxy razoável para responder ``\textit{quanto pesa cada
vertical na fatura mensal de storage}''.

\begin{table}[H]
\centering
\caption{\textbf{Custo \textit{storage} pro-rateado por \textit{vertical}},
        baseado na participação em \textit{bytes} de \textit{bronze}.}
\label{tab:storagepervert}
\begin{tabular}{lrrr}
\toprule
\textbf{Vertical}    & \textbf{Bronze (bytes)} & \textbf{Share \%} & \textbf{USD/mês} \\
\midrule
Bolsa Família        & 42,7 GB                 & 98,9              & 7,92 \\
Incontinência U.     & \phantom{0}488 MB       & \phantom{0}1,1    & 0,09 \\
Equipamentos médicos & \phantom{0}454 MB       & \phantom{0}1,0    & 0,08 \\
Emendas              & \phantom{00}3,7 MB      & \phantom{0}0,01   & 0,001 \\
\bottomrule
\end{tabular}
\end{table}

A leitura é evidente: \textbf{a fatura mensal de \textit{storage}
desta plataforma é, em essência, a fatura do \textit{vertical}
Bolsa Família}. Isso reflete a granularidade da fonte (CGU publica
microdados de pagamento mensal a beneficiário): manter histórico
2013\textendash 2025 íntegro custa \(\sim\)\us 7{,}92/mês.

% ---------- 4. DISCUSSÃO ------------------------------------------------
\section{Discussão}\label{sec:discussion}

\subsection{O padrão arquitetural emergente}

A vertical FinOps deste \textit{paper} adota um padrão simples e
reproduzível: \textit{system tables} do provedor \textit{cloud}
funcionam como \textit{bronze} ``virtual'' (não há \textit{ingest}
de dados externos \textemdash{} a Databricks já mantém o histórico
em sua própria infraestrutura), e os passos seguintes do
\textit{medallion} são equivalentes aos de qualquer outra
\textit{vertical}: \textit{silver} \textit{batch overwrite}
diário, \textit{gold} pivotado por \textit{persona} de consumidor,
\textit{export} JSON para \textit{front-end} estático.

Esse padrão tem três propriedades atraentes:

\begin{enumerate}\setlength{\itemsep}{4pt}
  \item \textbf{Auditoria como subproduto}. Não é necessária
        ferramenta externa de \textit{cost management}: a mesma
        infraestrutura que processa microdados públicos processa
        os \textit{records} de \textit{billing}. O custo marginal
        de operar a \textit{vertical} FinOps tende a zero
        \textemdash{} é mais um \textit{job} no \textit{bundle},
        com cadência diária e DBU consumido próximo do mínimo
        técnico.
  \item \textbf{Granularidade fina sem instrumentação adicional}.
        Cada \textit{job run} é registrado individualmente em
        \texttt{system.lakeflow.job\_run\_timeline}, sem
        necessidade de \textit{wrappers} de \textit{logging},
        \textit{exporters} OpenTelemetry, ou \textit{agents}
        proprietários instalados nos clusters.
  \item \textbf{Reprodutibilidade independente de fornecedor}. Tanto
        \textit{system tables} quanto \texttt{usage\_quantity} em
        DBU são padrões publicados pela Databricks; o algoritmo de
        \textit{pricing} versionado no tempo é trivialmente
        portado para AWS, Azure ou GCP. A vertical em si depende
        do fornecedor, mas o \textit{padrão} não.
\end{enumerate}

\subsection{Implicações da magnitude do \textit{wasted spend}}

A descoberta de que 53,9\,\% do gasto em \textit{jobs} consumiu
\textit{runs} fracassadas é desconfortavelmente próxima dos números
reportados pela literatura cinza FinOps em ambientes empresariais
de escala muito maior. FLEXERA (2024) reporta 32\,\% de
\textit{wasted spend} médio em organizações sem governança de custo
formalizada; o FinOps Foundation reporta entre 30\,\% e 35\,\% como
\textit{baseline}; este \textit{paper} reporta 54\,\% em ambiente
\textit{single-developer}.

Várias hipóteses parciais são consistentes com a magnitude
observada:

\begin{itemize}\setlength{\itemsep}{2pt}
  \item Em ambiente individual, há menos pressão de
        \textit{accountability}: nenhum \textit{stakeholder}
        cobra explicação por uma \textit{run} fracassada custando
        \us 3{,}01;
  \item O \textit{cost-per-run} desta plataforma (médio \us 0{,}40,
        máximo \us 3{,}01) é ordem de magnitude inferior ao
        \textit{cost-per-run} típico em produção empresarial,
        reduzindo o atrito psicológico para abortar/recomeçar;
  \item Erros de \textit{schema drift} em CSVs governamentais
        (CGU/CSV de Emendas Parlamentares com renomeações de coluna
        ano-a-ano; DATASUS \textit{layouts} variando entre
        \textit{competências}) produzem falhas idiossincráticas
        difíceis de antecipar em \textit{tests} sintéticos.
\end{itemize}

A \textbf{leitura geral} é que o ciclo
\textit{visibility} \textrightarrow{} \textit{allocation}
\textrightarrow{} \textit{optimization} parece estruturalmente
escala-invariante: mesmo uma plataforma com gasto mensal de
dezenas de dólares apresenta \textit{wasted spend} em proporção
similar à de plataformas empresariais. A ferramenta para o
combate \textemdash{} o \textit{ledger} desta auditoria
\textemdash{} tem custo marginal próximo de zero.

\subsection{Padrões aprendíveis para outros operadores}

Sem entrar em prescrição normativa, os números deste
\textit{paper} sugerem três focos para qualquer operador
individual de plataforma \textit{cloud}:

\begin{enumerate}\setlength{\itemsep}{2pt}
  \item Materializar uma vertical FinOps \textit{também} é um
        \textit{pipeline} \textit{open-source} reproduzível: ele
        custa USD~marginal próximo de zero e produz, gratuitamente,
        a contabilidade que de outro modo exigiria ferramentas
        \textit{third-party};
  \item As \textit{runs} \texttt{CANCELLED} de longa duração
        (média 50,6 min na nossa amostra) são candidatas naturais
        para \textit{circuit-breakers} (timeouts agressivos,
        \textit{budget caps} por \textit{job});
  \item \textit{Storage} domina a fatura recorrente para
        \textit{verticais} de microdados governamentais (Bolsa
        Família \(\sim\)99\,\% do nosso storage) \textemdash{}
        merece política explícita de retenção (\textit{vacuum},
        \textit{checkpoint} de tabelas históricas) antes que a
        cauda contínua de custo se torne dominante.
\end{enumerate}

% ---------- 5. LIMITAÇÕES -----------------------------------------------
\section{Limitações}\label{sec:limits}

\subsection{Janela de observação não-estacionária}

A janela \(d_0 = 11/06/2025\) a \(d_T = 28/04/2026\) é dominada
pelos últimos 20 dias. Estimadores construídos sobre
\(\bar{c}_{\text{lifetime}}\) (custo médio diário \textit{lifetime})
estão fortemente influenciados por essa concentração temporal e
não são bons preditores de gasto futuro. O \textit{run rate} de 30
dias é o estimador mais defensável para projeção; estimativas
\textit{lifetime} servem apenas como descrição histórica fiel.

\subsection{Latência das \textit{system tables}}

\texttt{system.billing.usage} apresenta latência aproximada de uma
hora; \textit{records} muito recentes podem não aparecer.
\texttt{job\_run\_timeline} pode demorar minutos a horas para
consolidar \textit{result\_state} final em \textit{runs}
canceladas. Resultado: o \textit{ledger} é eventualmente
consistente, e \textit{snapshots} produzidos a poucos minutos do
fim de uma \textit{run} podem subestimar o gasto correspondente.

\subsection{Atribuição de \textit{storage} ao \textit{vertical}}

A pro-rate de \texttt{DEFAULT\_STORAGE} pela participação em
\textit{bytes} de \textit{bronze} é aproximação \textemdash{} a
Databricks não emite \textit{records} de \textit{billing} com
\texttt{schema}/\texttt{table} no \texttt{usage\_metadata}. A
estimativa por \textit{vertical} é, portanto, contábil-gerencial,
não auditável diretamente contra a fatura.

\subsection{Free Edition vs.\ Standard/Premium}

Os preços por DBU em \texttt{system.billing.list\_prices} no Free
Edition refletem condições específicas do plano gratuito. Migração
da plataforma para um \textit{tier} pago (Standard, Premium) ou
para outro \textit{cloud} alteraria os \textit{list prices}, mas
não a metodologia: o \textit{join} versionado no tempo continua
válido com nova lista de preços. As magnitudes absolutas reportadas
neste \textit{paper} não são generalizáveis a outros \textit{tiers};
as \textit{proporções} (53,9\,\% \textit{wasted}, 7,3\,\%
\textit{storage}) são candidatas a verificação empírica.

\subsection{Plataforma \textit{single-developer}}

Esta auditoria documenta uma plataforma operada por um único
\textit{developer}. Plataformas com múltiplos
\textit{contributors} apresentam padrões adicionais
(\textit{cost-by-team}, \textit{chargeback}/\textit{showback})
não cobertos aqui. A questão sobre se a magnitude do
\textit{wasted spend} cresce, decresce ou permanece estável com
o tamanho da equipe permanece em aberto.

% ---------- 6. REPRODUTIBILIDADE ----------------------------------------
\section{Reprodutibilidade}\label{sec:reproducibility}

Este \textit{working paper} adota o seguinte protocolo de
reprodutibilidade:

\begin{itemize}\setlength{\itemsep}{2pt}
  \item Código-fonte do \textit{pipeline} (cinco \textit{notebooks}
        Python: 2 \textit{silver}, 2 \textit{gold}, 1 \textit{export};
        \(\sim\)600 linhas) disponível em
        \href{https://github.com/leonardochalhoub/mirante-dos-dados-br/tree/main/pipelines/notebooks}%
        {\texttt{github.com/leonardochalhoub/mirante-dos-dados-br/tree/main/pipelines/notebooks}}
        (acesso em 28/04/2026, 23h00 BRT).
  \item Configuração do \textit{Databricks Asset Bundle} (DAB)
        que orquestra o \textit{pipeline}, com \textit{schedule}
        diário 06h00 UTC, em \texttt{pipelines/databricks.yml}
        \texttt{job\_finops\_refresh}.
  \item \textit{Ledger} versionado em
        \texttt{data/gold/finops\_summary.json} no mesmo
        repositório \textemdash{} \textit{snapshot} reproduzível
        do estado das tabelas no momento da publicação.
  \item Metadata Unity Catalog (comentários por tabela e por
        coluna; \textit{tags} normalizadas com \texttt{layer},
        \texttt{domain}, \texttt{grain}, \texttt{currency},
        \texttt{refresh\_cadence}, \texttt{wp\_referenced}) aplicados
        idempotentemente via \texttt{job\_apply\_catalog\_metadata}.
  \item Algoritmo SQL canônico do \textit{pricing} versionado
        no tempo, com semântica
        \([\text{price\_start\_time}, \text{price\_end\_time})\),
        documentado na Seção~\ref{sec:methods} desta
        \textit{paper}.
\end{itemize}

A reprodução por terceiros em outro \textit{workspace} Databricks
exige apenas: (i) acesso a \texttt{system.billing.*} e
\texttt{system.lakeflow.*} (habilitados por padrão em Unity Catalog);
(ii) clone do repositório; (iii) deploy via
\texttt{databricks bundle deploy}. Não há \textit{secrets}, dados
proprietários, ou dependências externas além das próprias
\textit{system tables}.

% ---------- 7. CONCLUSÃO ------------------------------------------------
\section{Conclusão}

Este \textit{working paper} apresentou auditoria fim-a-fim do gasto
de uma plataforma de dados públicos brasileiros operada por um
único \textit{developer} no Databricks Free Edition, abrangendo a
janela completa de 322 dias entre 11 de junho de 2025 e 28 de
abril de 2026. A contribuição metodológica é a demonstração de que
\textit{system tables} podem servir como \textit{bronze} virtual
para uma \textit{vertical} FinOps integrada à mesma arquitetura
\textit{medallion} dos demais \textit{verticais} de microdados, com
\textit{silver} de granularidade fina e \textit{gold} pivotado para
consumo via \textit{front-end} estático. A contribuição empírica é
a documentação de \textit{wasted spend} de 53,9\,\% do gasto em
\textit{jobs}, magnitude próxima ou superior à reportada pela
literatura FinOps em ambientes empresariais de escala muito maior.
O ciclo
\textit{visibility} \textrightarrow{} \textit{allocation}
\textrightarrow{} \textit{optimization} parece estruturalmente
escala-invariante; o \textit{ledger} construído neste \textit{paper}
é a primeira metade desse ciclo, com custo marginal próximo de
zero.

% ---------- BIBLIOGRAFIA ------------------------------------------------
\section*{Referências}
\addcontentsline{toc}{section}{Referências}
\begin{singlespace}\small

\noindent ARMBRUST, M. et al. \textbf{Lakehouse: a new generation
of open platforms that unify data warehousing and advanced
analytics}. In: Conference on Innovative Data Systems Research
(CIDR), 11., 2021. \textit{Proceedings\ldots}\
2021. Disponível em:
\href{https://www.cidrdb.org/cidr2021/papers/cidr2021_paper17.pdf}%
{\texttt{cidrdb.org/cidr2021/papers/cidr2021\_paper17.pdf}}.
Acesso em: 28 abr.\ 2026, 23h00 BRT.

\vspace{6pt}
\noindent CHALHOUB, L. \textbf{Working Paper n.\ 2: Bolsa Família,
Auxílio Brasil e Novo Bolsa Família (2013\textendash 2025)}.
Brasil, 2026a. Disponível em:
\href{https://github.com/leonardochalhoub/mirante-dos-dados-br/blob/main/articles/bolsa-familia.tex}%
{\texttt{github.com/leonardochalhoub/mirante-dos-dados-br}}.
Acesso em: 28 abr.\ 2026, 23h05 BRT.

\vspace{6pt}
\noindent CHALHOUB, L. \textbf{Working Paper n.\ 7: Bolsa Família
ao nível municipal (5.570 unidades)}.
Brasil, 2026b. Disponível em:
\href{https://github.com/leonardochalhoub/mirante-dos-dados-br/blob/main/articles/bolsa-familia-municipios.tex}%
{\texttt{github.com/leonardochalhoub/mirante-dos-dados-br}}.
Acesso em: 28 abr.\ 2026, 23h05 BRT.

\vspace{6pt}
\noindent DATABRICKS. \textbf{System tables for billing,
auditing and pipeline observability}. Documentação técnica.
Disponível em:
\href{https://docs.databricks.com/en/admin/system-tables/index.html}%
{\texttt{docs.databricks.com/en/admin/system-tables}}.
Acesso em: 28 abr.\ 2026, 23h10 BRT.

\vspace{6pt}
\noindent DATABRICKS. \textbf{Unity Catalog: governance for
lakehouse}. Documentação técnica, 2024. Disponível em:
\href{https://docs.databricks.com/en/data-governance/unity-catalog/index.html}%
{\texttt{docs.databricks.com/en/data-governance/unity-catalog}}.
Acesso em: 28 abr.\ 2026, 23h12 BRT.

\vspace{6pt}
\noindent FINOPS FOUNDATION. \textbf{FinOps Framework}.
Documentação aberta, 2023. Disponível em:
\href{https://www.finops.org/framework/}%
{\texttt{finops.org/framework/}}.
Acesso em: 28 abr.\ 2026, 23h15 BRT.

\vspace{6pt}
\noindent FLEXERA. \textbf{State of the Cloud Report 2024}.
Itasca, IL: Flexera, 2024. Disponível em:
\href{https://info.flexera.com/CM-REPORT-State-of-the-Cloud}%
{\texttt{info.flexera.com/CM-REPORT-State-of-the-Cloud}}.
Acesso em: 28 abr.\ 2026, 23h18 BRT.

\vspace{6pt}
\noindent KIMBALL, R.; ROSS, M. \textbf{The Data Warehouse Toolkit:
the definitive guide to dimensional modeling}. 3.\ ed.\ Indianapolis:
John Wiley \& Sons, 2013.

\end{singlespace}

\end{document}
